


\section{Experimental Evaluation}
\label{sec:experiments}

We implemented a prototype of \system\ and evaluated it on several datasets prepared for the \emph{Text-to-Constraints} setting. Our evaluation addresses four questions: \textbf{RQ1}, how accurately can \system\ identify the sentences that express data constraints? \textbf{RQ2}, how accurately can it map them to executable DCs that capture the corresponding violations? \textbf{RQ3}, how do classifier and mapper implementations, LLM backbones, and prompt variants affect quality and robustness? \textbf{RQ4}, how does \system\ compare with data profiling and Text-to-SQL baselines?

\subsection{Datasets}
\label{sec:datasets}

%No benchmark currently exists for the \emph{Text-to-Constraints} problem.
%We prepared four datasets for our experiments. Two are based on real-world regulatory documents paired with corresponding tabular datasets, and two are derived from datasets widely used in data quality and profiling research, for which the target constraints are already known. For the real-world datasets, we manually annotated the constraints expressed in the source text. For the derived datasets, we generated natural-language descriptions of the known constraints, as described below. When a constraint had no natural violations in the data, we injected a few violating records so that every constraint could be evaluated under the violation-based metrics of Section~\ref{sec:metrics}. 
%Table~\ref{tab:datasets} summarizes the datasets.


We evaluate \system\ on six datasets: four based on real-world regulatory documents paired with corresponding tabular datasets, and two derived from benchmark datasets commonly used in data quality and profiling research. Table~\ref{tab:datasets} summarizes them. For every dataset, the authors labeled each sentence and wrote its gold DC in violation form over the released schema; sentences not checkable on the released attributes are non-constraints. No sentence occurs twice, and near-identical sentences (similarity $\geq 0.9$) are templated statements that differ in one attribute or literal (13 pairs in SISU, 32 in HMDA Checks, at most two elsewhere), never sharing a gold DC.



\begin{table}[ht]
  \centering
\small
%  \resizebox{\columnwidth}{!}{%
  \begin{tabular}{lrrrrr}
    \toprule
    \textbf{Dataset} & \textbf{\#Rows} & \textbf{\#Cols} & \textbf{Sent.} & \textbf{Const.} & \textbf{Viol.} \\
    \midrule
    Hospital & 1k & 17 & 109 & 15 & 834 \\
    Restaurant & 128k & 21 & 33 & 7 & 52,311 \\
    SISU & 102k & 56 & 42 & 17 & 533 \\
    Tax & 100k & 15 & 126 & 6 & 468 \\
    HMDAZ$^{\dagger}$ & 1M & 135 & 17 & 9 & 209 \\
    HMDA Checks$^{\dagger}$ & 1M & 81 & 158 & 92 & 986,491 \\
    \bottomrule
  \end{tabular}%
  %}% resizebox
    \caption{Dataset statistics. $^{\dagger}$First 1M rows of the public file (15.1M and 4.8M rows); \emph{Const.} counts sentences mapped to a non-null gold DC.}
  \label{tab:datasets}
\end{table}

\smallskip\noindent\textbf{SISU.}
This dataset contains admission records of Brazil's university selection system (Example~\ref{ex:intro}). We annotated three sections of \textit{Portaria Normativa No.~21/2012} covering student
registration, candidate classification and selection, and waitlist procedures.
The constraints span 11 nullability checks, a functional dependency, an order dependency (Example~\ref{ex:intro}), a conditional functional dependency (a student is selected in at most one option), and three conditional DCs over one or two tuples. Many have violations in the released data.


\smallskip\noindent\textbf{Restaurant.}
This dataset pairs the NYC inspection records (Example~\ref{ex:classification}) with the grading and closure sections of the food service inspection guide, which constrain score-to-grade mappings, inspection frequency, grade posting, and closure conditions. For the one rule over inspection sequences, we added a derived column with the days until the next inspection; automating such transformations is future work.


\smallskip\noindent\textbf{Tax and Hospital.}
Both datasets are standard benchmarks in DC research~\cite{chu2013discovering,holoclean}, but their constraints exist only in formal form. An LLM (Claude~4.5 Sonnet, not evaluated in any \system\ configuration) wrote a description of each constraint and surrounding non-constraint sentences in schema terminology, mimicking a document that interleaves rules with other prose. These two datasets nevertheless form a controlled setting: their constraints come from prior papers and their sentences from an LLM that may share priors with the evaluated backbones, which likely makes them easier than real regulatory text. SISU, Restaurant, and the HMDA datasets, which pair real regulatory text with real data, are the more reliable signal.




\smallskip\noindent\textbf{HMDAZ and HMDA Checks.}
Both datasets use public loan-level records released under the Home Mortgage Disclosure Act (HMDA); we use the first 1M rows of each file. \emph{HMDAZ} pairs the 2018 loan application
register (135 columns, including the code and label versions of each
categorical field) with Sections 1026.32, 1026.35, and 1026.43 of Regulation~Z,
which define high-cost (HOEPA) and higher-priced mortgages, the limitations
that apply to them, and the ability-to-repay requirements. Sentences are verbatim from
the regulation, and gold DCs keep the coverage conditions of each trigger.
Scope-only text, or text whose outcome is not observable in the data, is
annotated with a \texttt{null} constraint. \emph{HMDA Checks} pairs the 2014
register (81 columns) with the 2014 HMDA edit specification published by the
FFIEC: one entry per edit, with the edit test verbatim; edits that cannot be
checked on a single record (aggregate, transmittal-sheet, or external-reference
edits) are \texttt{null}. Compared with the other datasets, these two add wide schemas whose relevant column is often named after a code rather than the concept in the text, rules with arithmetic, and highly concentrated violations.

\subsection{Baselines}
\label{sec:baselines}

Since no prior system addresses the Text-to-Constraints task directly, we compare \system\ against the two closest paradigms: DC discovery from data alone and Text-to-SQL.


\smallskip\noindent\textbf{Data profiling.}
We use the DC discovery algorithms from~\cite{Pena2022FastDC}, running both exact and approximate denial-constraint discovery. Exact DCs hold with no violations and thus cannot match rules violated in the data, so we focus on approximate DCs, which tolerate a fraction $\varepsilon$ of violating pairs. Since FDs and order dependencies with at most two tuples are special cases of DCs, approximate DC discovery over the same attributes subsumes their discovery, and we do not run separate FD or OD algorithms. Conditional variants require predicates with constants, which DC discovery supports only for a restricted set of values; a comparison with dedicated CFD discovery is left for future work.

We run exact discovery and $\varepsilon \in \{0.001, 0.005, 0.01, 0.05, 0.1\}$, rank the DCs by coverage and succinctness~\cite{chu2013discovering}, and report, for each metric, the best result over this grid and top-$k \in \{5, 10, 20, 50, 100\}$.

Since discovery does not scale to the full SISU, Restaurant, and HMDA schemas, we restrict them to the columns of the gold constraints, the most favorable setting for the baseline, and run it on a seeded sample of 100k of the 1M HMDA rows. On HMDA Checks (39 columns), no configuration finished: runs timed out or exhausted a 12\,GB heap even with samples of 10k rows.


\smallskip\noindent\textbf{Text-to-SQL.}
We prompt XiYan-SQL~\cite{XiYanSQL}, a state-of-the-art Text-to-SQL framework, to return for each sentence a query that selects the violating tuples, or \texttt{NULL} if the sentence is not a checkable constraint, so it acts as both classifier and mapper. We use the 32B model with three seeds and report the mean.


\subsection{Evaluation Metrics}
\label{sec:metrics}

We evaluate three dimensions; as in Text-to-SQL, execution-based evaluation complements structural matching, because exact form alone does not capture correctness~\cite{zhong-etal-2020-semantic}.
\emph{Constraint classification} measures precision, recall, and F$_1$ of the sentences labeled as constraints; it applies to \system\ and Text-to-SQL, not to profiling.
\emph{Constraint mapping} (CM) compares predicted and gold DCs by exact predicate-set match (precision, recall, F$_1$). For Text-to-SQL, predicates are read from the \texttt{WHERE} clause (a correlated \texttt{EXISTS} as a pair DC); a query with no DC reading, e.g., with aggregation or \texttt{CASE}, counts as correct if it returns exactly the violations of the gold DC.
\emph{Violation detection} executes all predicted and all gold constraints and compares the sets of tuples that take part in at least one violation (aggregate violation F$_1$, AV); it applies to all methods.

\smallskip\noindent\textbf{Complementary mapping metrics.}
Exact matching penalizes logically equivalent rewritings, and the tuple-level aggregate can be dominated by a single rule. We therefore also report:
(i)~\emph{equivalence-aware mapping} (CM$_{\equiv}$), which compares DCs after a logical normal form (e.g., \texttt{IN} as a disjunction, sorted operands) and unifies literals on a code column and its label column (e.g., \texttt{hoepa\_status = 1} and \texttt{hoepa\_status\_name = 'High-cost mortgage'}) only when the data confirms that they select the same tuples;
(ii)~\emph{execution equivalence}, the fraction of gold rules with violations whose predicted DC flags exactly the same tuples;
(iii)~\emph{per-rule violation F$_1$}, averaged over gold sentences so that every rule weighs the same; and
(iv)~\emph{violation concentration}, the share of gold-violating tuples flagged by the largest gold rule. Predicted DCs whose SQL fails are counted as errors, not as producing no violation.

\subsection{Implementation and Settings}
\label{sec:implementation}

\system\ is implemented in Python with DuckDB as the execution backend. All LLM calls go through OpenRouter with strict JSON-schema outputs. We evaluate six backbones: GPT-6 Luna, DeepSeek-V4-Flash, Mistral Large 4, Qwen3.8-2.4T-A95B, GPT-6 Sol, and Claude Sonnet 5.5. Two candidates were replaced: Kimi K2.6, which emitted whitespace until the token limit under strict decoding, and MiniMax-M3, which returned empty single-prompt outputs (Section~\ref{sec:constraintextraction}). Classifier and mapper configurations run once, with up to 32k output tokens per call; end-to-end configurations run three times. A preliminary study with earlier backbones (e.g., Qwen3-32B and GPT-5) followed the same protocol on Hospital, Restaurant, SISU, and Tax; we refer to it in Section~\ref{sec:constraintextraction}. \system\ and DC discovery ran on a workstation with an Intel Core i7-7700HQ CPU and 16\,GB of memory, and XiYan-SQL on an NVIDIA A100 (80\,GB). The baselines are scored by the same evaluator, data, and query limits as \system.
